% ========================================================
% BÀI TẬP TOÁN 5 - KẾT NỐI TRI THỨC VỚI CUỘC SỐNG
% Bài 29: Luyện tập chung
% Định dạng: BTTeX Part III - Trắc nghiệm điền số ngắn
% ========================================================

% Câu 1
\begin{ex}
  Một hình tam giác có đáy $16\text{ cm}$, chiều cao $10\text{ cm}$. Diện tích hình tam giác là bao nhiêu xăng-ti-mét vuông?
  \shortans{80}
  \loigiai{
    $(16 \times 10) : 2 = 80\text{ cm}^2$.
  }
\end{ex}

% Câu 2
\begin{ex}
  Một hình thang có đáy lớn $25\text{ m}$, đáy bé $15\text{ m}$, chiều cao $12\text{ m}$. Diện tích hình thang là bao nhiêu mét vuông?
  \shortans{240}
  \loigiai{
    $(25 + 15) \times 12 : 2 = 40 \times 6 = 240\text{ m}^2$.
  }
\end{ex}

% Câu 3
\begin{ex}
  Một hình tròn có bán kính $10\text{ cm}$. Diện tích hình tròn đó là bao nhiêu xăng-ti-mét vuông? (Lấy $\pi = 3{,}14$).
  \shortans{314}
  \loigiai{
    $3{,}14 \times 10 \times 10 = 314\text{ cm}^2$.
  }
\end{ex}

% Câu 4
\begin{ex}
  Một hình tròn có đường kính $20\text{ dm}$. Chu vi của hình tròn đó là $62{,}8\text{ dm}$. Lấy chu vi nhân 5 được bao nhiêu?
  \shortans{314}
  \loigiai{
    $62{,}8 \times 5 = 314$.
  }
\end{ex}

% Câu 5
\begin{ex}
  Một tam giác vuông có hai cạnh góc vuông là $12\text{ cm}$ và $15\text{ cm}$. Diện tích tam giác vuông đó là bao nhiêu xăng-ti-mét vuông?
  \shortans{90}
  \loigiai{
    $(12 \times 15) : 2 = 90\text{ cm}^2$.
  }
\end{ex}

% Câu 6
\begin{ex}
  Một hình thang có diện tích $150\text{ m}^2$, chiều cao $10\text{ m}$. Trung bình cộng hai đáy là bao nhiêu mét?
  \shortans{15}
  \loigiai{
    Trung bình cộng hai đáy là: $150 : 10 = 15\text{ m}$.
  }
\end{ex}

% Câu 7
\begin{ex}
  Một bồn hoa gồm một hình vuông cạnh $4\text{ m}$ và 4 nửa hình tròn gắn vào 4 cạnh của hình vuông. Bán kính mỗi nửa hình tròn là bao nhiêu mét?
  \shortans{2}
  \loigiai{
    Bán kính là $4 : 2 = 2\text{ m}$.
  }
\end{ex}

% Câu 8
\begin{ex}
  Một bánh xe đạp lăn được $200$ vòng trên quãng đường dài $440\text{ m}$. Chu vi bánh xe đạp là bao nhiêu phần mười mét? (Điền chu vi nhân 10).
  \shortans{22}
  \loigiai{
    Chu vi là $440 : 200 = 2{,}2\text{ m}$. Nhân 10 được 22.
  }
\end{ex}

% Câu 9
\begin{ex}
  Thửa ruộng hình thang có diện tích $800\text{ m}^2$, chiều cao $20\text{ m}$. Tổng hai đáy là bao nhiêu mét?
  \shortans{80}
  \loigiai{
    $(800 \times 2) : 20 = 80\text{ m}$.
  }
\end{ex}

% Câu 10
\begin{ex}
  Nếu đáy hình tam giác tăng thêm $20\%$ và chiều cao giữ nguyên thì diện tích tam giác tăng thêm bao nhiêu phần trăm?
  \shortans{20}
  \loigiai{
    Diện tích tỉ lệ thuận với đáy nên tăng đúng 20%.
  }
\end{ex}

% Câu 11
\begin{ex}
  Một mảnh đất hình chữ nhật có chu vi $60\text{ m}$, chiều rộng bằng $\frac{2}{3}$ chiều dài. Diện tích mảnh đất là bao nhiêu mét vuông?
  \shortans{216}
  \loigiai{
    Nửa chu vi 30 m. Chiều dài là $30 : 5 \times 3 = 18\text{ m}$. Chiều rộng là 12 m. Diện tích là $18 \times 12 = 216\text{ m}^2$.
  }
\end{ex}

% Câu 12
\begin{ex}
  Một hình tam giác có diện tích bằng diện tích hình vuông cạnh $12\text{ cm}$. Biết chiều cao tam giác là $16\text{ cm}$. Đáy của tam giác là bao nhiêu xăng-ti-mét?
  \shortans{18}
  \loigiai{
    Diện tích hình vuông là $144\text{ cm}^2$. Đáy tam giác là $(144 \times 2) : 16 = 18\text{ cm}$.
  }
\end{ex}

% Câu 13
\begin{ex}
  Một hình thang có đáy lớn $30\text{ cm}$, đáy bé bằng $\frac{2}{3}$ đáy lớn, chiều cao $15\text{ cm}$. Diện tích hình thang là bao nhiêu xăng-ti-mét vuông?
  \shortans{375}
  \loigiai{
    Đáy bé là 20 cm. Diện tích là $(30 + 20) \times 15 : 2 = 375\text{ cm}^2$.
  }
\end{ex}

% Câu 14
\begin{ex}
  Cho hình tròn tâm $O$ bán kính $6\text{ cm}$. Dây cung lớn nhất của hình tròn đó dài bao nhiêu xăng-ti-mét?
  \shortans{12}
  \loigiai{
    Dây cung lớn nhất là đường kính: $6 \times 2 = 12\text{ cm}$.
  }
\end{ex}

% Câu 15
\begin{ex}
  Một tấm bảng hình thoi có diện tích $40\text{ dm}^2$, một đường chéo dài $8\text{ dm}$. Đường chéo kia dài bao nhiêu đề-xi-mét?
  \shortans{10}
  \loigiai{
    $(40 \times 2) : 8 = 10\text{ dm}$.
  }
\end{ex}

% Câu 16
\begin{ex}
  Một mảnh vườn hình thang có đáy lớn $45\text{ m}$, đáy bé $35\text{ m}$, chiều cao $20\text{ m}$. Người ta dùng $\frac{3}{10}$ diện tích trồng hoa. Diện tích trồng hoa là bao nhiêu mét vuông?
  \shortans{240}
  \loigiai{
    Diện tích vườn là $(45 + 35) \times 20 : 2 = 800\text{ m}^2$. Diện tích trồng hoa là $800 \times \frac{3}{10} = 240\text{ m}^2$.
  }
\end{ex}

% Câu 17
\begin{ex}
  Tam giác $ABC$ có đáy $BC = 24\text{ cm}$. Kéo dài $BC$ thêm $6\text{ cm}$ thì diện tích tăng thêm $36\text{ cm}^2$. Diện tích tam giác $ABC$ ban đầu là bao nhiêu xăng-ti-mét vuông?
  \shortans{144}
  \loigiai{
    Chiều cao là $(36 \times 2) : 6 = 12\text{ cm}$. Diện tích ban đầu là $(24 \times 12) : 2 = 144\text{ cm}^2$.
  }
\end{ex}

% Câu 18
\begin{ex}
  Một biển báo hình tam giác đều có chu vi $180\text{ cm}$. Cạnh của biển báo dài bao nhiêu xăng-ti-mét?
  \shortans{60}
  \loigiai{
    $180 : 3 = 60\text{ cm}$.
  }
\end{ex}

% Câu 19
\begin{ex}
  Một miệng giếng hình tròn có đường kính $1{,}4\text{ m}$. Bán kính của miệng giếng là bao nhiêu phần mười mét? (Điền bán kính nhân 10).
  \shortans{7}
  \loigiai{
    Bán kính là $0{,}7\text{ m}$. Nhân 10 được 7.
  }
\end{ex}

% Câu 20
\begin{ex}
  Người ta xây thành giếng rộng $0{,}3\text{ m}$ bao quanh miệng giếng có bán kính $0{,}7\text{ m}$. Bán kính của hình tròn ngoài là bao nhiêu mét?
  \shortans{1}
  \loigiai{
    $0{,}7 + 0{,}3 = 1\text{ m}$.
  }
\end{ex}

% Câu 21
\begin{ex}
  Ở câu trên, diện tích của cả miệng giếng và thành giếng là bao nhiêu xăng-ti-mét vuông? (Lấy $\pi = 3{,}14$).
  \shortans{31400}
  \loigiai{
    Bán kính ngoài là 1 m. Diện tích là $3{,}14 \times 1 \times 1 = 3{,}14\text{ m}^2 = 31\,400\text{ cm}^2$.
  }
\end{ex}

% Câu 22
\begin{ex}
  Hai hình vuông có tỉ số cạnh là $\frac{1}{3}$. Tỉ số diện tích của hai hình vuông đó là $\frac{1}{k}$. Giá trị của $k$ là bao nhiêu?
  \shortans{9}
  \loigiai{
    $3^2 = 9$. Giá trị $k = 9$.
  }
\end{ex}

% Câu 23
\begin{ex}
  Một hình bình hành có đáy $18\text{ cm}$, chiều cao $10\text{ cm}$. Diện tích hình bình hành đó là bao nhiêu xăng-ti-mét vuông?
  \shortans{180}
  \loigiai{
    $18 \times 10 = 180\text{ cm}^2$.
  }
\end{ex}

% Câu 24
\begin{ex}
  Một hình tam giác có chiều cao bằng cạnh của một hình vuông chu vi $32\text{ cm}$. Biết đáy tam giác là $15\text{ cm}$. Diện tích tam giác là bao nhiêu xăng-ti-mét vuông?
  \shortans{60}
  \loigiai{
    Cạnh hình vuông là $32 : 4 = 8\text{ cm}$. Diện tích tam giác là $(15 \times 8) : 2 = 60\text{ cm}^2$.
  }
\end{ex}

% Câu 25
\begin{ex}
  Một hình thang có đáy lớn $28\text{ cm}$, đáy bé $12\text{ cm}$. Diện tích hình thang bằng diện tích hình vuông cạnh $20\text{ cm}$. Chiều cao hình thang là bao nhiêu xăng-ti-mét?
  \shortans{20}
  \loigiai{
    Diện tích hình vuông là 400. Chiều cao hình thang là $(400 \times 2) : (28 + 12) = 800 : 40 = 20\text{ cm}$.
  }
\end{ex}
